% ============================================================
%  CAPÍTULO 1 — INTRODUCCIÓN Y CONTEXTO
% ============================================================
\chapter{Introducción y contexto}

\section{Descripción del encargo}

El presente informe ha sido elaborado por el Departamento de Ingeniería de una empresa
especializada en construcción y mantenimiento de plantas industriales, en el marco de un
contrato de mantenimiento y explotación de una Estación Depuradora de Aguas Residuales
(EDAR) de titularidad pública.

El contrato establece la necesidad de estudiar y valorar determinadas mejoras de
infraestructura que contribuyan a una mejor operatividad y mantenimiento de las
instalaciones. Conforme a las instrucciones del Director General, el presente documento se
centra exclusivamente en la parte de \textbf{Obra Civil}, abarcando los dos apartados
indicados en el \textit{Anexo I} del contrato:

\begin{enumerate}
  \item \textbf{Muro de contención / pantalla}: contención del desnivel necesario para la
        implantación de una planta fotovoltaica, con una longitud en planta de 150~ml y un
        desnivel de \alturaDesnivelMuro~m.
  \item \textbf{Arqueta de fangos mezcla}: nuevo depósito de almacenamiento de fangos de
        dimensiones $\largoArqueta$~m~$\times$~$\anchoArqueta$~m~$\times$~$\alturaArqueta$~m
        (largo~$\times$~ancho~$\times$~alto).
\end{enumerate}

\section{Datos de partida}

\subsection{Parámetros del grupo}

Los parámetros $X$ e $Y$ se han obtenido a partir de la media de los DNIs de los seis
integrantes del grupo (tabla~\ref{tab:dnis}):

\begin{table}[H]
  \centering
  \caption{DNIs del grupo y cálculo de la media}
  \label{tab:dnis}
  \begin{tabular}{@{}cc@{}}
    \toprule
    \textbf{Alumno} & \textbf{DNI (número)} \\
    \midrule
    1 & 30\,248\,832 \\
    2 & 49\,505\,322 \\
    3 & 77\,936\,142 \\
    4 & 01\,647\,207 \\
    5 & 49\,275\,542 \\
    6 & 53\,965\,819 \\
    \midrule
    \textbf{Suma} & 262\,578\,864 \\
    \textbf{Media} & 43\,763\,144 \\
    \midrule
    $X$ & \paramX \\
    $Y$ & \paramY \\
    \bottomrule
  \end{tabular}
\end{table}

\subsection{Condicionantes especiales}

\begin{notabox}[Condicionantes medioambientales y geotécnicos]
  \begin{itemize}
    \item La zona de ejecución está \textbf{especialmente protegida medioambientalmente},
          por lo que se deben adoptar medidas específicas contra posibles filtraciones de
          carácter químico o bacteriano al terreno.
    \item Los análisis del estudio geotécnico indican una \textbf{concentración de
          sulfatos} de $3\,500$~mg~SO$_4^{2-}$/kg de suelo seco, lo que condiciona la
          elección del tipo de cemento y clase de exposición del hormigón.
  \end{itemize}
\end{notabox}

\subsection{Columna estratigráfica — Sondeo 1}

La tabla~\ref{tab:sondeo1} recoge los parámetros geotécnicos de los estratos identificados
en el Sondeo~1 (Reactor Biológico), que servirá de referencia para todos los cálculos
geotécnicos de este informe.

\begin{table}[H]
  \centering
  \caption{Columna estratigráfica Sondeo 1}
  \label{tab:sondeo1}
  \resizebox{\textwidth}{!}{%
  \begin{tabular}{@{}lcccccccc@{}}
    \toprule
    \textbf{Estrato} & \textbf{Cota (m)} &
    $\gamma$ \textbf{(Tn/m³)} & $\phi'$ \textbf{(°)} &
    $N_{\text{SPT}}$ & $C'$ \textbf{(kg/cm²)} &
    $C_u$ \textbf{(kg/cm²)} & $E$ \textbf{(MPa)} & $\nu$ \\
    \midrule
    Arena limosa gris-verdosa, poco densa
      & 0{,}00 a $-$2{,}50 & 1{,}82 & 29 & 9  & 0    & 0    & 6{,}894  & 0{,}25 \\
    \quad (saturada)
      & $-$2{,}50 a $-$3{,}50 & 1{,}98 & 29 & 9  & 0    & 0    & 6{,}894  & 0{,}25 \\
    Arcillas versicolores + intercalac. arena, consistencia media
      & $-$3{,}50 a $-$8{,}00 & 1{,}78 & 25{,}2 & 8  & 0{,}47 & 1{,}00 & 12{,}000 & 0{,}30 \\
    Arcillas margosas, consistencia firme
      & $-$8{,}00 a $-$25{,}00 & 1{,}92 & 26 & 33 & 0{,}78 & 1{,}52 & 25{,}278 & 0{,}25 \\
    \bottomrule
  \end{tabular}}
\end{table}

\begin{notabox}[Nivel freático]
  El nivel freático (N.F.) se sitúa a $-$2{,}50~m de profundidad, en el contacto con
  la capa de arena limosa saturada.
\end{notabox}
